\section{Results on competitive programming}\label{sec:results2}
\subsection{Textbook moves}
\begin{table}[t]\centering\small
\caption{Textbook moves added to the DGM seed (Haiku~5.5; all 232 fresh problems, two runs each, one interleaved batch).
E: string-replace editor; B: non-interactive shell; L: model loop with a 16{,}000-token budget; P: workflow prompt.
Gain over the seed in the same batch with 95\% intervals; cost in dollars per thousand tasks; share of runs that end
with no solution.}\label{tab:grid}
\begin{tabular}{lrrrr}\toprule
moves & $F$ (\%) & gain [95\% CI] & cost & no solution (\%) \\ \midrule
\csname @@input\endcsname tab_grid.tex
\bottomrule\end{tabular}\end{table}
Every combination that contains the model-loop move L gains between $\vgridWithLmin$ and $\vgridWithLmax$ points over
the seed, and every combination without it between $\vgridNoLmin$ and $\vgridNoLmax$ (Table~\ref{tab:grid}). L bundles
four changes (a 16{,}000-token budget, an answer to every tool call, one re-prompt after a truncated output, and a
40-turn limit); alone it gains $\vgridLgain$ points. The editor and the shell change nothing measurable on their own, and the workflow prompt alone costs
$\vgridPgainAbs$ points: its longer reasoning runs into the seed's 4{,}096-token cap more often.

\paragraph{The registered test (H6).} The second pre-registration asks whether the best textbook combination comes
within three points of the best first-round loop agent, the one from loop 3. The best of the 15 combinations is chosen
on $F_1$ and compared on $F_2$ (\S\ref{sec:round2}); we also named two single candidates to report on all of $F$, all
four moves and L. Loop 3's agent was re-evaluated in another batch, so these contrasts cross batches. We report raw
contrasts here and seed-adjusted ones in Appendix~\ref{app:batches}. All contrasts in this paragraph are the candidate minus loop 3's agent, in points.

The combination chosen on $F_1$ (\vHsixWinner) is $\vHsixFtwoAbs$ points [$\vHsixFtwoAbslo$, $\vHsixFtwoAbshi$] behind
on $F_2$, so H6 is not supported. Over the \vlcN{} combinations that contain L, the contrasts on $F_2$ range from
$\vlcFtwoMin$ to $\vlcFtwoMax$, and the chosen one happens to be the lowest. On all of $F$, all four moves are
$\vpsEBLPRawAbs$ points [$\vpsEBLPRawAbslo$, $\vpsEBLPRawAbshi$] behind and L is $\vpsLRawAbs$ [$\vpsLRawAbslo$,
$\vpsLRawAbshi$] behind. Of the \vlcN{} L combinations, \vlcMiss{} miss the margin and \vlcSig{} are significantly behind.
The literal registered reading chooses the best combination on all of $F$ (\vhsixLitBest) and scores it on the same
problems, which favours it; its contrast, $\vhsixLit$ [$\vhsixLitlo$, $\vhsixLithi$], is within the margin on raw scores
but not after seed adjustment. Every combination is cheaper per task than loop 3's agent, at most
$\vgridMaxCostShare\%$ of its cost; L alone costs $\vgridLCostShare\%$ as much, and all four moves
cost $\vgridEBLPCostShare\%$. Measured from the seed, the textbook moves thus close about nine tenths of the distance to the
best loop agent and remain one to five points behind it: on all fresh problems the contrasts of the L combinations range
from $\vlcFMax$ to $\vlcFMin$, and on $F_2$ the chosen combination is at $\vHsixFtwo$.

\paragraph{All loops against one move (exploratory).} The final batch measured every loop's reported best agent, the
seed, the cap-only seed and L together, two runs per problem (Table~\ref{tab:loops}). There L gains $\vlvLgain$ points
[$\vlvLgainlo$, $\vlvLgainhi$] over the seed. No loop agent is significantly better than L and $\vlvWorse$ of $\vlvN$
are significantly worse; the best, loop 3's agent, is $\vlvBest$ points [$\vlvBestlo$, $\vlvBesthi$] ahead, so
non-inferiority of L within three points is not shown in this batch either. In the third-round batch, loop 3's agent is
$\rcLoopThreeLAbs$ points [$\rcLoopThreeLAbslo$, $\rcLoopThreeLAbshi$] ahead of L, a significant gap
(\S\ref{sec:random-clean}). Section~\ref{sec:round1} reports how far the cap-only seed trails loop 3's agent in the final
and ablation batches. The proxy caps every request at 16{,}000 output tokens, so loop 3's agent never receives more output
than L; what it adds is more attempts after an empty turn. With one more attempt when its first solution fails the
sample tests in the problem statement (the rule of the compute-matched baseline above), L reaches $\vLtwoAcc\%$ at
\$$\vLtwoUsd$ per thousand tasks, while loop 3's agent reaches $\vLthreeAcc\%$ at \$$\vLthreeUsd$, a contrast of
$\vltwoRaw$ points [$\vltwoRawlo$, $\vltwoRawhi$]. This comparison pairs L's runs from the grid batch with loop 3's agent
from the ablation batch (Appendix~\ref{app:batches}); it was specified after the H6 result and before it was computed.

\subsection{One-shot redesigns}
\begin{table}[t]\centering\small
\caption{One-shot redesigns of the DGM seed by a language model (Haiku~5.5 task model; fresh problems, two runs each).
Each redesign saw the seed's code and five failure logs once. Cap: the output cap the redesign set. Opus redesigns are
registered; Sonnet redesigns are exploratory.}\label{tab:redesign}
\begin{tabular}{lrrrr}\toprule
agent & cap & $F$ (\%) & gain [95\% CI] & cost \\ \midrule
\csname @@input\endcsname tab_redesign.tex
\bottomrule\end{tabular}\end{table}
Four of the five Opus redesigns and three of the five Sonnet redesigns raise the output cap to 8{,}192 tokens and gain
18 to 22 points; the three that keep 4{,}096 change the pass rate by $-1$ to $+8$ points (Table~\ref{tab:redesign}). For
the registered H8 we take contrasts as redesign minus loop, again across batches (Appendix~\ref{app:batches}). The
contrast of the average Opus redesign with the average loop best is $\vHeightOpusmean$ [$\vHeightOpusmeanlo$,
$\vHeightOpusmeanhi$], so this part is supported. The best of five, chosen on $F_1$, is $\vHeightOpusbestAbs$ points
[$\vHeightOpusbestAbslo$, $\vHeightOpusbestAbshi$] behind loop 3's agent on $F_2$, so that part is not. The best Opus
redesign on all of $F$ is $\vheightLitOpusRawAbs$ points [$\vheightLitOpusRawAbslo$, $\vheightLitOpusRawAbshi$] behind.

\paragraph{Without failure logs.}\label{sec:redesign-logfree} In the third round, five further samples from Claude Opus~5.5 received the redesign
prompt without the five failure logs: the agent summary, the seed's code and the same task paragraph. Four of the five
raise the output cap to 8{,}192 tokens, as four of five did with the logs; the fifth keeps 4{,}096 but continues after a
truncated response and reruns the agent when no solution was written. In one batch with the seed, which scores $\lfSeedF\%$, the
log-free redesigns gain $\lfGainMin$ to $\lfGainMax$ points, on average $\lfMeanFree$ [$\lfMeanFreelo$, $\lfMeanFreehi$],
compared with $\lfMeanBased$ [$\lfMeanBasedlo$, $\lfMeanBasedhi$] for the redesigns made from failure logs; the difference is
$\lfDiff$ points [$\lfDifflo$, $\lfDiffhi$], so the registered non-inferiority test passes. An exploratory comparison shows that the
averages agree for different reasons. Among the redesigns that raise the cap, those made without logs gain $\lfRaiseDiffAbs$
points [$\lfRaiseDiffAbslo$, $\lfRaiseDiffAbshi$] less than those made from logs. Of the two that keep 4{,}096, the
log-free one gains $\lfKeepFree$ points and the log-based one gains $\lfKeepBased$. Claude Opus~5.5 thus finds the fix from the code
alone when asked for a one-shot redesign, so the failure logs are not needed to find it. In the first round, none of the
\nladConeatt{} diagnoses that Claude Sonnet~5.5 made from the agent's code alone mentions the output cap, stop reasons or
empty responses (Appendix~\ref{app:round1}).


\subsection{Search without diagnosis, another search rule, more loops}
\begin{table}[t]\centering\small
\caption{Reported best agents of all Haiku~5.5 loops, re-evaluated together with the seed in one batch (all 232 fresh
problems, two runs each). $S$: the selection score for which the loop reported the agent. Inflation: selection gain
over the seed's pooled selection score ($\nseedDgmSmean\%$ over $\nseedDgmSn$ evaluations) minus fresh gain. Cost in dollars per thousand tasks. The first-round DGM loops and the second-round
arms (more DGM loops, random diagnosis, HGM-style) are registered; the loops without hint prompts are exploratory. The
last rows are references measured in the same batch.}\label{tab:loops}
\begin{tabular}{llrrrr}\toprule
loop type & loop & $S$ & fresh gain [95\% CI] & inflation & cost \\ \midrule
\csname @@input\endcsname tab_loops.tex
\bottomrule\end{tabular}\end{table}
The random-diagnosis loops draw every request from the fixed list of \S\ref{sec:round2}; because one request targets
empty responses, they search over textbook requests rather than without any diagnosis. In the final batch they gain $\vhRandom$ points [$\vhRandomlo$, $\vhRandomhi$] on average,
whereas the first-round DGM loops gain $\vhDgmTwo$ [$\vhDgmTwolo$, $\vhDgmTwohi$]. The registered non-inferiority test
(H7) fails: the difference is $\vhRandomMinusDgm$ [$\vhRandomMinusDgmlo$, $\vhRandomMinusDgmhi$]. Against all six DGM
loops it is $\vhRandomMinusDgmAll$ [$\vhRandomMinusDgmAlllo$, $\vhRandomMinusDgmAllhi$]. HGM-style parent selection gains
$\vhHgm$ [$\vhHgmlo$, $\vhHgmhi$], which is $\vhHgmMinusDgmAllAbs$ [$\vhHgmMinusDgmAllAbslo$, $\vhHgmMinusDgmAllAbshi$]
points more than the six DGM loops (Table~\ref{tab:loops}). These intervals resample problems (\S\ref{sec:round2}). Compared at the
level of loops (three or six per arm), none of the arm differences is significant (Welch tests: HGM-style vs DGM
$p=\vllHgmP$, random vs first-round DGM $p=\vllRandomThreeP$, random vs all DGM $p=\vllRandomSixP$). Four of the 14 loops
reach the level of the single move L: loop 3, both loops without hint prompts and one HGM-style loop; one of the six DGM
loops does.

For each DGM loop the selection gain exceeds the fresh gain, by $\vhtwoMinInfl$ to $\vhtwoMaxInfl$ points. With each
loop's own sampling error (tasks resampled, as registered for H2) every interval includes zero, and over the six loops
the hierarchical estimate is $\vhtwoMean$ points [$\vhtwoLo$, $\vhtwoHi$]. The placebo loops' best selection scores
are given in \S\ref{sec:round1}.


\paragraph{With a clean request list.}\label{sec:random-clean}
Five loops ran the DGM procedure with every diagnosis replaced by a request drawn from a clean list, and three DGM loops
ran at the same time. The clean list is the second-round list with the three failure-related requests described in
\S\ref{sec:round2} replaced by three generic ones. The one aimed at empty responses also
asked the agent to answer every tool call, which is part of move L. In the third-round batch, which also re-evaluated the earlier loops' agents and where the seed scores $\rcSeedF\%$, the five clean
random-diagnosis loops
gain $\rcClean$ points [$\rcCleanlo$, $\rcCleanhi$] on average, with individual loops from $\rcCleanMin$ to
$\rcCleanMax$. All nine DGM loops gain $\rcDgmAll$ [$\rcDgmAlllo$, $\rcDgmAllhi$], and the three that ran at the same
time $\rcDgmNew$ [$\rcDgmNewlo$, $\rcDgmNewhi$]. The registered non-inferiority test fails: the clean loops trail the DGM loops
by $\rcHfourteenAbs$ points [$\rcHfourteenAbslo$, $\rcHfourteenAbshi$], and at the level of loops the difference is
significant (Welch test, $p=\rcHfourteenP$). The registered secondary contrast, against the three same-time DGM loops,
is $\rcCleanNew$ [$\rcCleanNewlo$, $\rcCleanNewhi$] ($p=\rcCleanNewP$).

The remaining contrasts in this batch are exploratory. The random-diagnosis loops of the second round gain $\rcRandOld$
[$\rcRandOldlo$, $\rcRandOldhi$]. That is
$\rcCleanOldAbs$ points [$\rcCleanOldAbslo$, $\rcCleanOldAbshi$] more than the clean loops ($p=\rcCleanOldP$) and level
with all nine DGM loops ($\rcRandAll$ [$\rcRandAlllo$, $\rcRandAllhi$]). The same-time DGM loops reproduce the earlier
ones within $\rcNewOldAbs$ points. Search over generic requests therefore gains much less when no request points at the
failure, and the DGM's diagnosis, which reads the failure logs, does better. The single model-loop move L gains $\rcL$
[$\rcLlo$, $\rcLhi$] in this batch, more than the average DGM loop but $\rcLoopThreeLAbs$ points [$\rcLoopThreeLAbslo$,
$\rcLoopThreeLAbshi$] less than the agent of loop 3.


\subsection{What the lineages contain}
The loops' edits depend on each other. Removing one edit from a lineage and re-applying the rest leaves rejected hunks
for all but \vlinThreeClean{} of loop 3's \vlinThreeN{} edits (\vlinTwoClean{} of \vlinTwoN{} for loop 2, \vlinOneClean{}
of \vlinOneN{} for loop 1), so the registered test of whether two edits carry most of the gain (H9) cannot be run.
Removing loop 3's last edit, the only leave-one-out variant that applies cleanly, costs $\vlasteditAbs$ points
[$\vlasteditAbslo$, $\vlasteditAbshi$]. In an exploratory variant, setting the output cap of loop 3's agent back to 4{,}096 tokens while keeping
its other code costs only $\vsemcapdropAbs$ points [$\vsemcapdropAbslo$, $\vsemcapdropAbshi$], probably because the
agent also retries empty turns with nudges and fallback calls.
